\section*{अध्याय \ref{ch:b1:stratigraphic-scale} --- स्तरिक पैमाना}

\begin{solution}{exo:b1:stratigraphic-scale:1}
(क)~कालस्तरिक, एक \omterm{def:b1:stratigraphic-scale:chronostratigraphy}{तंत्र}। (ख)~अश्मस्तरिक, एक \omterm{def:b1:stratigraphic-scale:lithostratigraphy}{संरचना}।
(ग)~भूकालिक, एक युग (उसकी चट्टान निम्न क्रिटेशियस \omterm{def:b1:stratigraphic-scale:chronostratigraphy}{श्रेणी} है)।
(घ)~जैव-स्तरिक, एक \omterm{def:b1:stratigraphic-scale:biozone}{जैव-क्षेत्र}। (ङ)~कालस्तरिक, एक \omterm{def:b1:stratigraphic-scale:chronostratigraphy}{चरण}।
(च)~भूकालिक, एक युग। (छ)~अश्मस्तरिक, एक \omterm{def:b1:stratigraphic-scale:lithostratigraphy}{सदस्य}।
\end{solution}

\begin{solution}{exo:b1:stratigraphic-scale:2}
पैलियोज़ोइक $538.8 - 251.902 = \qty{286.9}{\Myr}$, फ़ैनेरोज़ोइक का \qty{53.2}{\%};
मेसोज़ोइक $251.902 - 66.00 = \qty{185.9}{\Myr}$, \qty{34.5}{\%}; सेनोज़ोइक
\qty{66.0}{\Myr}, \qty{12.2}{\%}। फ़ैनेरोज़ोइक $538.8/4567 = \qty{11.8}{\%}$ है
हेडियन के आधार से समय का।
\end{solution}

\begin{solution}{exo:b1:stratigraphic-scale:3}
A, B, C, D उसी क्रम में जमे (अध्यारोपण), क्षैतिज रूप से (मूल क्षैतिजता), इसलिए
उनका झुकाव एक बाद की घटना है। डाइक E, A से D को काटती है, इसलिए वह D से नई है
(आर-पार संबंध); झुकाव से पहले या बाद में, यह विवरण से नहीं बताया जा सकता। अपरदन-पृष्ठ
D और E को काटता है, इसलिए दोनों से नया है (आर-पार संबंध)। F और G उस पर बाद में
जमे (अध्यारोपण)। भ्रंश A से F तक विस्थापित करता है पर G को नहीं: वह F के बाद और G से
पहले चला (आर-पार संबंध)। क्रम: A, B, C, D; झुकाव और E का अंतर्वेधन (क्रम अनिर्णीत);
अपरदन; F; भ्रंशन; G।
\end{solution}

\begin{solution}{exo:b1:stratigraphic-scale:4}
$r = 30/(85.40 - 84.60) = \qty{37.5}{m}$ प्रति दस लाख वर्ष। पहली उपस्थिति निचली राख
से $12/37.5 = \qty{0.32}{\Myr}$ बाद है: $85.40 - 0.32 =
\qty{85.08}{\Ma}$।
\end{solution}

\begin{solution}{exo:b1:stratigraphic-scale:5}
(क)~शैल-स्तंभ के विरुद्ध ऊँचाई पर चार ऊर्ध्वाधर पट्टियाँ। (ख)~B और C \qtyrange{26}{50}{m}
पर साथ मिलते हैं: एक समवर्ती-परास क्षेत्र \qty{24}{m} मोटा। B की पहली उपस्थिति से
D की पहली उपस्थिति तक का \omterm{def:b1:stratigraphic-scale:biozone}{अंतराल-क्षेत्र} \qtyrange{12}{41}{m} के बीच चलता है, \qty{29}{m}।
(ग)~दूसरे खंड में B (\qty{3}{m}) और C (\qty{9}{m}) की पहली उपस्थितियाँ रखी जा सकती
हैं, इसलिए (ख) के \omterm{def:b1:stratigraphic-scale:biozone}{अंतराल-क्षेत्र} का आधार और समवर्ती-परास क्षेत्र का आधार पहचाने
जाते हैं; उसका शीर्ष (D की पहली उपस्थिति) और खंड 1 में दिखे अतिव्यापन का शीर्ष जाँचे
नहीं जा सकते, क्योंकि D अनुपस्थित है। B और C की पहली उपस्थितियाँ ($\qty{12}{m} \leftrightarrow
\qty{3}{m}$, $\qty{26}{m} \leftrightarrow \qty{9}{m}$) सबसे सुरक्षित रेखाएँ हैं; अंतिम उपस्थितियाँ
(B \qty{50}{m} और \qty{15}{m} पर) संरक्षण पर अधिक निर्भर करती हैं।
\end{solution}

\begin{solution}{exo:b1:stratigraphic-scale:6}
(क)~$60/20 = \qty{3}{\Myr}$: भूमि की ओर वाले कुएँ में शीर्ष बलुआ पत्थर तीन दस लाख
वर्ष छोटा है। (ख)~एक दस लाख वर्ष की कीचड़, \qty{30}{m} प्रति दस लाख वर्ष: राख
बलुआ पत्थर के शीर्ष से \qty{30}{m} ऊपर है। (ग)~राख गिरने के समय तट-रेखा समुद्र की
ओर वाले कुएँ से \qty{20}{km} भूमि की ओर खिसक चुकी थी, इसलिए वह भूमि की ओर वाले कुएँ
से अब भी \qty{40}{km} समुद्र की ओर थी, जो सूखी भूमि पर खड़ा था। राख वहाँ भूमि की
सतह पर गिरी, उस बलुआ पत्थर के नीचे जिसे समुद्र दो दस लाख वर्ष बाद जमाता है ---
यदि वह पहले बह न गई हो, जैसा भूमि पर एक ढीली परत के साथ आसानी से हो सकता है। वही
राख-परत एक कुएँ में संरचना-सीमा के ऊपर और दूसरे में उसके नीचे है: सीमा समय को पार
करती है।
\end{solution}

\begin{solution}{exo:b1:stratigraphic-scale:7}
(क)~लोचकोविअन $1.36/419.62 = \qty{0.32}{\%}$; इंडुअन $0.024/251.902 =
\qty{0.0095}{\%}$,
लगभग 34 गुना छोटा। इंडुअन का आधार \omterm{def:b1:stratigraphic-scale:gssp}{GSSP} खंड में ही सटीक आँकी गई राख-परतों के बीच पड़ता
है; एक बड़ी अनिश्चितता बताती है कि यह एक ऐसी सीमा है जिसकी आयु दूर के आँकड़ों से
ले जानी होगी। (ख)~$251.941 - 251.880 = \qty{0.061}{\Myr} = 61$ हज़ार वर्ष, अनिश्चितता
$\sqrt{37^2 + 31^2} = 48$ हज़ार वर्ष के साथ: प्रकाशित
\num[uncertainty-mode=separate]{60 \pm 48} हज़ार वर्ष।
\end{solution}

\begin{solution}{exo:b1:stratigraphic-scale:8}
(क)~$m = 3000 \times \frac43\pi\,(5000)^3 = \qty{1.57e15}{kg}$; धूल, $60m = \qty{9.42e16}{kg}$।
(ख)~पृथ्वी की सतह $4\pi(6.378\times10^6)^2
= \qty{5.11e14}{m^2}$: प्रति वर्ग मीटर \qty{184}{kg}। (ग)~$184/2500 =
\qty{0.074}{m}$, लगभग \qty{7}{cm}। (घ)~मोटाई व्यास के घन के अनुपात में बढ़ती है:
\qty{6}{km} के लिए \qty{1.6}{cm}, \qty{14}{km} के लिए \qty{20}{cm}। परिमाण की कोटि,
सेंटीमीटर से डेसीमीटर तक, सीमा-मृत्तिका की है, जो एल केफ़ पर \qty{50}{cm} मोटी है;
तुलना केवल मोटे तौर पर हो सकती है, क्योंकि धूल समान रूप से नहीं गिरी होगी, और
मृत्तिका अपने समय का सामान्य तलछट भी बटोरती है।
\end{solution}

\begin{solution}{exo:b1:stratigraphic-scale:9}
(क)~$\num{60e3} \times \num{2e-6} = \qty{0.12}{m}$, \qty{12}{cm}; 12 और 108 हज़ार वर्षों
के लिए \qty{2.4}{cm} और \qty{21.6}{cm}। (ख)~$\num{32e3} \times
\num{10e-6} = \qty{0.32}{m}$। (ग)~एक विलोपन चट्टान के सेंटीमीटर से डेसीमीटर तक घेरता
है: नमूने सेंटीमीटर-दर-सेंटीमीटर लेने होंगे, और संघनित खंड में \qty{10}{cm} मोटा एक
नमूना पचास हज़ार वर्षों का औसत देता है --- घटना का अधिकांश भाग।
\end{solution}

\begin{solution}{exo:b1:stratigraphic-scale:10}
(क)~$u = R/s = 10 \times 2000 = \qty{20000}{m}$ प्रति दस लाख वर्ष,
\qty{20}{km} प्रति दस लाख वर्ष; \qty{30}{km} पर, $\Delta t = \qty{1.5}{\Myr}$।
(ख)~मैदान अब सतह $z = s x + a t$ है। तट-रेखा $s x + a t = z_0 + R t$ को संतुष्ट करती
है, इसलिए $x(t) = (z_0 + (R - a)t)/s$ और $\mathrm dx/\mathrm dt = (R - a)/s$, जो तब तक
धनात्मक (भूमि की ओर) है जब तक $a < R$ है। जब $a > R$, तो गति ऋणात्मक है: तट-रेखा
समुद्र की ओर खिसकती है, एक पश्चगमन, और तट का बलुआ पत्थर तब समुद्र की ओर युवा है।
(ग)~\qty{200}{m} लंबी एक चट्टान-दीवार, \qty{20}{km} प्रति दस लाख वर्ष पर, $0.2/20 = \qty{0.01}{\Myr}$ फैलती है --- दस हज़ार वर्षों का तिर्यक-कालिकता, जो जीवाश्म या अधिकांश
आँकड़े जितना विभेदित कर सकते हैं उससे कहीं कम; सीमा के दोनों ओर की चट्टान उसके साथ-साथ
हर जगह एक ही है।
\end{solution}

\begin{solution}{exo:b1:stratigraphic-scale:11}
(क)~$f = (25 - 10)/(50 - 10) = 0.375$ और $t = 100.50 - 0.375 \times 0.80 =
\qty{100.20}{\Ma}$, के साथ
\[
  \sigma_t = \sqrt{(0.625 \times 0.05)^2 + (0.375 \times 0.03)^2} = \qty{0.033}{\Myr}.
\] (ख)~$g(f) = (1 - f)^2\sigma_1^2 + f^2\sigma_2^2$ का $g'(f) = -2(1 - f)\sigma_1^2 + 2f\sigma_2^2$ है, जो $f^* =
\sigma_1^2/(\sigma_1^2 + \sigma_2^2) = 0.0025/0.0034 = 0.735$ पर शून्य है, बेहतर आँकी गई परत के निकट; वहाँ $\sigma_t = \sigma_1\sigma_2/\sqrt{\sigma_1^2 + \sigma_2^2} =
\qty{0.026}{\Myr}$, दोनों में से किसी से छोटा। (ग)~$g$ ऊपर खुलने वाली एक परवलय है, इसलिए $[0, 1]$ पर वह एक सिरे पर सबसे बड़ा है: $g(0) = \sigma_1^2$, $g(1) = \sigma_2^2$। अतः दोनों परतों के बीच $\sigma_t \le \max(\sigma_1, \sigma_2)$ है।
\end{solution}

\begin{solution}{exo:b1:stratigraphic-scale:12}
(क)~परत की गहराई में आधी पर सीमा है: $f = 0.5$, \qty{29.50}{\Ma}। (ख)~$\varphi_0\lambda = 0.63 \times 1960 = \qty{1234.8}{m}$ के साथ, \qtyrange{600}{1000}{m} से ठोस मोटाई $400 - 1234.8\,(e^{-600/1960} - e^{-1000/1960}) =
\qty{232.2}{m}$ है, और \qtyrange{200}{600}{m} से \qty{194.2}{m}; कुल मिलाकर \qty{426.3}{m}। निचली राख और सीमा के बीच जमे कणों का अंश $232.2/426.3 = 0.545$ है, और आयु $30.0 - 0.545 =
\qty{29.46}{\Ma}$ है। (ग)~दूसरा अनुमान \num{45} हज़ार वर्ष छोटा है। वह बेहतर है: गहरी आधी में पानी अधिक खो चुका है, इसलिए वह प्रति मीटर उथली आधी से अधिक कण, और अधिक समय, रखती है (औसत छिद्रता \num{0.42} बनाम \num{0.51})।
\end{solution}

\subsection*{समस्या: एक द्रोणी के पार तीन खंडों का सहसंबंध}

\begin{solution}{pb:b1:stratigraphic-scale:1}
\begin{enumerate}
\item तीन स्तंभ, मान लें \qty{10}{m} के लिए \qty{1}{cm} पर, राख-परतों और क्षेत्र-सीमाओं से जुड़े; केल बलुआ पत्थर का पच्चर A में \qty{45}{m} से C में \qty{6}{m} तक पतला होता है।
  \item क्षेत्र~2: A $65.8 - 36.6 = \qty{29.2}{m}$, B \qty{20.8}{m}, C \qty{10.0}{m}। क्षेत्र~3 लॉग किया गया: A \qty{14.2}{m}, B \qty{14.0}{m}, C \qty{7.2}{m}।
  \item क्षेत्र~1 \emph{वार्डेनोसेरास प्राइमम} की पहली उपस्थिति पर शुरू होता है, जो हर खंड के आधार से मिलता है: उसकी पहली उपस्थिति उन सबके नीचे है।
  \item वह केवल बलुआ पत्थर में रहता है: वह संस्तर-प्रकार का अनुसरण करता है, और संस्तर-प्रकार स्थान-स्थान पर अलग-अलग आयु का है। A में वह क्षेत्र~1 और~2 में मिलता है, B और C में केवल क्षेत्र~1 में।
  \item तीनों में क्षेत्र~2 में (A \qty{49.0}{m}, B \qty{40.0}{m}, C \qty{19.0}{m})। यहाँ वह खंड-दर-खंड सुसंगत है, और इसलिए क्षेत्र~2 को उप-विभाजित कर सकता है, पर अंतिम उपस्थिति संरक्षण पर निर्भर करती है; उसे पहले और खंडों में जाँचना चाहिए।
  \item रेखाएँ \qty{36.6}{m}, \qty{31.2}{m} और \qty{14.8}{m} को जोड़ती हैं, फिर \qty{65.8}{m}, \qty{52.0}{m} और \qty{24.8}{m}। क्षेत्र~2 $29.2 : 20.8 : 10.0$ के अनुपात में पतला होता है, लगभग $2.9 : 2.1 : 1$: यदि क्षेत्र हर जगह समान समय तक चला, तो तलछट A में C की तुलना में लगभग तीन गुना तेज़ी से जमा।
  \item A: क्षेत्र~2 ($45.0 > 36.6$); B: क्षेत्र~1 ($20.0 < 31.2$); C: क्षेत्र~1 ($6.0 < 14.8$)।
  \item नहीं: वह A में क्षेत्र~2 में, B और C में क्षेत्र~1 में है, इसलिए वह पश्चिम में युवा है। समुद्री शेल तट-बालू के ऊपर पश्चिम की ओर फैला: तट-रेखा पश्चिम की ओर, भूमि की ओर, एक ट्रांसग्रेशन में खिसकी।
  \item एक राख पूरी द्रोणी पर एक भूवैज्ञानिक क्षण में गिरती है और किसी भी तलछट पर बैठती है; बलुआ पत्थर का शीर्ष खिसकती तट-रेखा का अनुसरण करता है और तिर्यक-कालिक है।
  \item पार्श्व सातत्य (परत द्रोणी के आर-पार जारी रहती है), और अध्यारोपण (राख~1 हर खंड में राख~2 के नीचे, जो सहसंबंध की जाँच करता है)।
  \item A $(71.0 - 5.0)/1.20 = \qty{55.0}{m}$ प्रति दस लाख वर्ष; B $48.0/1.20 = \qty{40.0}{}$; C $24.0/1.20 = \qty{20.0}{}$।
  \item A भूमि के सबसे निकट है, जहाँ तलछट की आपूर्ति थी; पूर्व की गहरी जल तक उसका कम भाग पहुँचा।
  \item $0.1/55 = \num{1.8}$ हज़ार वर्ष A में, $\num{2.5}$ B में, $\num{5.0}$ C में: A समय को सबसे बारीकी से विभेदित करता है।
  \item A $152.40 - 40.0/55.0 = \qty{151.67}{\Ma}$; B $152.40 - 12.0/40.0 =
        \qty{152.10}{\Ma}$; C $152.40 - 3.0/20.0 = \qty{152.25}{\Ma}$।
  \item C से B: $20/(152.25 - 152.10) = \qty{133}{km}$ प्रति दस लाख वर्ष; B से A: $25/(152.10 - 151.67) = \qty{58.5}{km}$ प्रति दस लाख वर्ष। वह धीमा हुआ।
  \item A $152.40 - 60.8/55.0 = \qty{151.29}{\Ma}$; B $152.40 - 44.0/40.0 =
        \qty{151.30}{\Ma}$; C $152.40 - 21.8/20.0 = \qty{151.31}{\Ma}$।
  \item वे \num{16} हज़ार वर्षों में फैलते हैं, जो हर आयु की अनिश्चितता (लगभग \num{28} हज़ार वर्ष, प्रश्न~19) से कम है: पहली उपस्थिति तिथियों के विभेदन के भीतर पूरी द्रोणी में समकालिक है, और यही उसे एक अच्छा सूचक बनाता है।
  \item A $152.40 - 31.6/55.0 = \qty{151.826}{\Ma}$; B $152.40 - 23.2/40.0 =
        \qty{151.820}{\Ma}$; C $152.40 - 11.8/20.0 = \qty{151.810}{\Ma}$।
  \item $f = 0.479$, $0.483$, $0.492$; $\sigma_t = \sqrt{(1 - f)^2 + f^2} \times
        0.04 = \qty{0.028}{\Myr}$ हर एक में, \num{28} हज़ार वर्ष।
  \item माध्य \qty{151.819}{\Ma}; फैलाव $151.826 - 151.810 = \num{16}$ हज़ार वर्ष, हर आयु की अनिश्चितता से छोटा: तीनों सहमत हैं।
  \item प्रक्षेपण राख~1 से राख~2 तक एक ही दर मानता है। A में उस अंतराल में \qty{45}{m} तक तट-बलुआ पत्थर और ऊपर शेल है; रेत और कीचड़ एक ही दर से जमे होने ज़रूरी नहीं, और वे अलग तरह से संपीड़ित होते हैं (\cref{ch:b1:diagenesis-basins})। यदि रेत तेज़ी से जमी, तो बलुआ पत्थर वाला भाग उसकी मोटाई से सुझाए गए से कम समय रखता है, और पहली उपस्थिति परिकलित से पुरानी है --- A से आयु वास्तव में तीनों में सबसे पुरानी है।
  \item B: पहली उपस्थिति समुद्री शेल में है, बलुआ पत्थर के शीर्ष से \qty{11}{m} ऊपर, स्तर के निकट संस्तर-प्रकार में बिना परिवर्तन के, और दर \qty{10}{cm} पर \num{2.5} हज़ार वर्ष देती है। A समय को बेहतर विभेदित करता है, पर स्तर को तट-बलुआ पत्थर में रखता है, संस्तर-प्रकार के परिवर्तन से \qty{8.4}{m} नीचे; C संघनित है, \qty{10}{cm} पर \num{5} हज़ार वर्ष। विभेदन के लिए C सबसे खराब है, संस्तर-प्रकार के लिए A।
  \item दो पहली उपस्थितियों की माध्य आयुओं से, $151.819 - 151.302 =
        \qty{0.52}{\Myr}$: लगभग आधा दस लाख वर्ष।
  \item सीमा है \qty[uncertainty-mode=separate]{151.82(3)}{\Ma}।
\end{enumerate}
\end{solution}
